\section*{\texorpdfstring{\S\CurrentExerciseGroup}{Section \CurrentExerciseGroup}}
\phantomsection
\addcontentsline{toc}{section}{\texorpdfstring{Exercises for \S\CurrentExerciseGroup}{Exercises for Section \CurrentExerciseGroup}}

\begin{enumerate}
\def\labelenumi{\arabic{enumi})}
\item
  Show that, for a numerical function \(f\) on \(T\) to be essentially integrable for every positive measure on \(T\) , it is necessary and sufficient that it be bounded, have compact support, and be measurable for every positive measure on \(T\) .
\item
  Let \(G\) be a barreled space, \(H\) a semi-reflexive space, and \(F\) the space \(\mathcal{L}_s(G;H)\) . Let \(m\) be a vectorial measure on \(T\) , with values in \(F\) . Show that for every numerical function \(f\) essentially integrable for \(m\) , \(\int fd\mathbf{m} \in F\) . (Observe that \(F\) is semi-reflexive.) The case where \(H = R\) .
\item
  Let \(\mathbf{m}\) be a vectorial measure on \(\mathrm{T}\) , with values in \(\mathrm{F}\) . Show that for every numerical function \(f\) essentially integrable for \(\mathbf{m}\) , one has \(\int f dm \in \mathrm{F}''\) in each of the following two cases: \(a)\) \(\mathrm{F}\) is distinguished (TVS, IV, §2, Exer. 4); \(b)\) \(\mathrm{F}\) is a Fréchet space and \(\mathrm{T}\) is countable at infinity. (In both cases, embed \(\mathrm{F}\) in \(\mathrm{F}''\) and apply the method of Cor. 1 of Prop. 3; in case \(b\) ), use the Cor. of Prop. 3 of TVS, IV, §3, No.~3.)
\item
  Let \(\mathbf{m}\) be a vectorial measure on \(\mathrm{T}\) , with values in \(\mathrm{F}\) . Show that for every numerical function \(f \geqslant 0\) that is lower semi-continuous and essentially integrable for \(\mathbf{m}\) , one has \(\int f d\mathbf{m} \in \mathrm{F}''\) (use Th. 1 of Ch. IV, §1, No.~1).
\item
  Let F be a semi-reflexive locally convex space, m a measure on T with values in F. Show that for every numerical function g locally integrable for m, the mapping \(f \mapsto \int fg dm\) of \(\mathcal{K}(T)\) into F is a measure on T, denoted \(g \cdot m\) ; in order that a numerical function h be essentially integrable for \(g \cdot m\) , it is necessary and sufficient that gh be essentially integrable for m, in which case \(\int h d(g \cdot m) = \int hg dm\) .
\item
  Let F be a semi-reflexive locally convex space, equipped with an algebra structure over R, such that the mapping \((u,v)\mapsto uv\) of \(F\times F\) into F is separately continuous. Let m be a measure on T, with values in F, such that \(\mathbf{m}(fg)=\mathbf{m}(f)\mathbf{m}(g)\) for f and g in \(\mathcal{K}(\mathrm{T})\) . Show that if the numerical functions f,g are such that f,g, and fg belong to \(\mathcal{L}(\mathbf{m})\) , then \(\int fg dm=(\int fd\mathbf{m})(\int gd\mathbf{m})\) . (First treat the case that \(g\in\mathcal{K}(\mathrm{T})\) , by considering the two measures \(f\mapsto\mathbf{m}(fg)\) and \(f\mapsto\mathbf{m}(f)\mathbf{m}(g)\) ; pass to the general case by proceeding similarly and making use of Exer. 5.) The case where \(F=\mathcal{L}(\mathrm{G})\) , G being reflexive and F equipped with the topology of pointwise convergence or the topology of weak pointwise convergence.
\item
  Let \(\mu\) be a positive measure on T, m a measure on T with values in F, having base \(\mu\) and density f with respect to \(\mu\) . Let \(q\) be a lower semi-continuous semi-norm on F. Show that if, for every compact subset K of T, \(\int_{K}^{*}(q \circ f) d\mu\) is finite, then, for every \(\mu\) -measurable function \(h \geqslant 0\) , the inequality \(\int^{\bullet} h dq(m) \leqslant \int^{\bullet} (q \circ f) h d\mu\) holds. (Use Lemma 1 of Ch. V, 1st edn., §2, No.~2.)²
\item
  Let \(\mathbf{m}\) be a vectorial measure on \(\mathbf{T}\) with values in \(\mathbf{F}\) , \(q\) -majorizable for a lower semi-continuous semi-norm \(q\) on \(\mathbf{F}\) . Show that for every function \(f \geqslant 0\) of \(\mathcal{H}(\mathbf{T})\) ,
\end{enumerate}

\[
\int f d (q \circ \mathbf {m}) = \sup \sum_ {i} q (\mathbf {m} (f _ {i})),
\]

where \((f_i)\) runs over the set of all finite sequences of elements of \(\mathcal{H}(\mathrm{T})\) such that \(\sum_{i}|f_i| \leqslant f\) . (Argue as in Th. 1 of Ch. II, §2, No.~2.)

\begin{enumerate}
\def\labelenumi{\arabic{enumi})}
\setcounter{enumi}{8}
\tightlist
\item
  Let T be a locally compact space countable at infinity, F a Hausdorff barreled space containing a countable dense subset, m a vectorial measure on T with values in the weak dual F' of F. Show that there exists a positive measure \(\mu\) on T such that m is scalarly of base \(\mu\) . (Make use of Prop. 12 of Ch. V, §5, No.~6, and condition 5) of Cor. 5 of the Lebesgue--Nikodym theorem, Ch. V, §5, No.~5, Th. 2.)
\end{enumerate}

¶ 10) Let K be a compact space. For every Hausdorff quotient space \(K_{1}\) of K, one identifies the Banach space \(\mathcal{C}(K_{1})\) with a closed subspace of \(\mathcal{C}(K)\) by means of the injection \(f \mapsto f \circ \pi\) , where \(\pi\) is the canonical mapping of K onto \(K_{1}\) .

\begin{enumerate}
\def\labelenumi{\alph{enumi})}
\item
  Let \(u\) be a continuous linear mapping of \(\mathcal{C}(\mathrm{K})\) into a quasi-complete Hausdorff locally convex space \(\mathrm{F}\) . In order that \(u\) transform the unit ball of \(\mathcal{C}(\mathrm{K})\) into a subset of \(\mathrm{F}\) that is relatively compact for \(\sigma(\mathrm{F},\mathrm{F}')\) , it suffices that, for every metrizable quotient space \(\mathrm{K}_1\) of \(\mathrm{K}\) , the restriction of \(u\) to \(\mathcal{C}(\mathrm{K}_1)\) have that property. (Make use of Eberlein's theorem (TVS, IV, §5, No.~3, Th. 1), and observe that every sequence \((f_n)\) in \(\mathcal{C}(\mathrm{K})\) is contained in a subspace \(\mathcal{C}(\mathrm{K}_1)\) , where \(\mathrm{K}_1\) is a metrizable quotient space of \(\mathrm{K}\) ; for that, one considers the continuous mapping \(x \mapsto (f_n(x))\) of \(\mathrm{K}\) into \(\mathbf{R}^{\mathbf{N}}\) .)
\item
  For every Hausdorff quotient space \(K_{1}\) of K, the transpose of the canonical injection of \(\mathcal{C}(K_{1})\) into \(\mathcal{C}(K)\) is the mapping \(\mu\mapsto\pi(\mu)\) of \(\mathcal{M}(K)\) into \(\mathcal{M}(K_{1})\) . Denote by \((\mathcal{M}(\mathrm{K}))'\) the dual of the Banach space \(\mathcal{M}(\mathrm{K})\) (the bidual of \(\mathcal{C}(\mathrm{K})\) ). For a subset A of \(\mathcal{M}(\mathrm{K})\) to be relatively compact for \(\sigma(\mathcal{M}(\mathrm{K}),(\mathcal{M}(\mathrm{K}))')\) , it is necessary and sufficient that, for every metrizable quotient space \(K_{1}\) of K, the canonical image of A in \(\mathcal{M}(\mathrm{K}_{1})\) be relatively compact for \(\sigma(\mathcal{M}(\mathrm{K}_{1}),(\mathcal{M}(\mathrm{K}_{1}))')\) . (Observe that A is strongly bounded; one can then suppose that A is convex, balanced and closed for the vague topology \(\sigma(\mathcal{M}(\mathrm{K}),\mathcal{C}(\mathrm{K}))\) (apply in \(\mathcal{M}(\mathrm{K}_{1})\) Exer. 10 of Ch. IV, §7); A is therefore the image of the unit ball of the dual \(F'\) of a Banach space F, under the transpose of a continuous linear mapping u of \(\mathcal{C}(\mathrm{K})\) into F (consider the polar of A in \(\mathcal{C}(\mathrm{K})\) ); then apply a).)
\end{enumerate}

¶ 11) Let F, G be two Hausdorff locally convex spaces, G being assumed quasi-complete; let u be a continuous linear mapping of F into G; its transpose \(^{t}u\) being strongly continuous, the transpose \(^{t}(^{t}u)\) is a strongly and weakly continuous linear mapping of F'\,' into G'\,'. Show that the following conditions are equivalent:

\(\alpha)\) \(u\) transforms every bounded subset of \(\mathbf{F}\) into a relatively compact subset of \(\mathbf{G}\) for \(\sigma(\mathbf{G},\mathbf{G}')\) .

\(\beta)^{t}(^{t}u)\) maps \(\mathbf{F}^{\prime \prime}\) into G.

\(\gamma)\ ^{t}u\) is continuous for the topologies \(\sigma(\mathrm{G}^{\prime},\mathrm{G})\) and \(\sigma(\mathrm{F}^{\prime},\mathrm{F}^{\prime\prime})\) .

\(\delta)^{t}u\) transforms every equicontinuous subset of \(\mathrm{G}'\) into a set relatively compact for \(\sigma (\mathbf{F}',\mathbf{F}'')\) .

(Use the fact that every bounded subset of F is relatively compact in \(F''\) for \(\sigma(F'', F')\) . To see that \(\delta\) implies \(\gamma\) ), consider first the case that G is complete, and use Cor. 1 of Th. 2 of TVS, III, §3, No.~6. To pass to the general case, embed G in its completion and observe that every point of \(F''\) is in the closure, for the topology \(\sigma(F'', F')\) , of a bounded subset of F.)

¶ 12) a) Let K be a compact space, E = E(K) the subspace of R \(^{K}\) formed by the linear combinations of characteristic functions of open sets; one identifies E with a subspace of the bidual of C(K). Show that a Cauchy sequence for σ(M(K),E) is convergent for σ(M(K),(M(K))') (make use of Prop. 12 and Exer. 15 of Ch. V, §5). From this, deduce that every subset A of M(K) that is relatively compact for the topology σ(M(K),E), is relatively compact for the topology σ(M(K),(M(K))'). (Reduce to the case that K is metrizable, using Exer. 10 b). Show that A is bounded, by applying Exer. 13 of Ch. V, §5. Note that on A, the vague topology σ(M(K),C(K)) and the topology σ(M(K), \(\overline{E}\) ) are identical, \(\overline{E} \supset C(K)\) being the closure of E in the strong dual of M(K); deduce from this, using the fact that C(K) is separable, that every sequence of points of A has a subsequence that is Cauchy for σ(M(K),E); conclude by invoking Eberlein's theorem.)

\begin{enumerate}
\def\labelenumi{\alph{enumi})}
\setcounter{enumi}{1}
\item
  Let T be a locally compact space, m a vectorial measure on T with values in a quasi-complete, Hausdorff locally convex space F. Assume that, for every compact subset K of T, \(\int_{K} dm \in F\) ; show that, for every bounded Borel f function on T with compact support, one has \(\int f dm \in F\) , and that the image under \(f \mapsto \int f dm\) of the set of Borel functions with support contained in a compact subset K of T and of norm \(\|f\| \leqslant 1\) , is relatively compact for \(\sigma(F, F')\) (make use of a) and Exer. 11, applied to \(u : f \mapsto \int f dm\) , where f runs over the set of linear combinations of characteristic functions of compact subsets of T).
\item
  Let F be a quasi-complete, Hausdorff locally convex space, m a vectorial measure on T with values in F, scalarly of base \(\mu\) ; for every \(z^{\prime} \in F^{\prime}\) , write \(z^{\prime} \circ m = g_{z^{\prime}} \cdot \mu\) . For the condition of b) to be fulfilled, it is necessary and sufficient that, for every compact subset K of T, every equicontinuous subset \(H^{\prime}\) of \(F^{\prime}\) , and every \(\varepsilon > 0\) , there exist a \(\delta > 0\) such that the relations \(A \subset K\) and \(\mu^{*}(A) \leqslant \delta\) imply \(\int_{A}^{*}|g_{z^{\prime}}| d\mu \leqslant \varepsilon\) for every \(z^{\prime} \in H^{\prime}\) (use Exer. 11 above and Exer. 15 of Ch. V, §5); m is then said to be absolutely continuous with respect to \(\mu\) (for the original topology of F). Every vectorial measure with values in a semi-reflexive space and scalarly of base \(\mu\) , is absolutely continuous with respect to \(\mu\) (Cor. 2 of Prop. 3). Every majorizable vectorial measure m with values in a Banach space \(\mathbf{F}\) is absolutely continuous with respect to \(|\mathbf{m}|\) (for the original topology of \(\mathbf{F}\) ).
\end{enumerate}

¶ 13) Let G be a Hausdorff barreled space, F = G' its dual; if, when F is equipped with a topology compatible with the duality between F and G, m is a vectorial measure on T with values in F, then m is again a vectorial measure when F is equipped with the strong topology of the dual of G.

Assume m to be scalarly of base \(\mu\) ; then m is absolutely continuous with respect to \(\mu\) , when F is equipped with the topology \(\tau(\mathrm{F},\mathrm{G})\) (Exer. 12 c)). In order that m be absolutely continuous with respect to \(\mu\) when F is equipped with the strong topology, it is necessary and sufficient that, for every compact subset K of T and every sequence \((\mathrm{A}_{n})\) of \(\mu\) -measurable subsets of K, the image B under m of the closure in \(\mathcal{L}^{1}(\mu)\) of the set of step functions of absolute value \(\leqslant1\) , over the clan generated by the sequence of the \(A_{n}\) (Ch. IV, §4, No.~9), contain a dense countable subset for the strong topology of \(G'\) . (To see that the condition is sufficient, reduce to the case that T = K is a Stone space (Ch. IV, §4, Exer. 10). Then argue by contradiction: if \((\mathrm{A}_{n})\) is a sequence of subsets of K that are both open and closed, one can assume, on passing to a metrizable quotient space \(K_{1}\) of K (Exer. 10 a)), that the continuous functions on \(K_{1}\) corresponding to the \(\varphi_{A_{n}}\) form a total set in \(\mathcal{C}(\mathrm{K}_{1})\) . On the other hand, on considering the quotient of G by the subspace orthogonal to B, one can suppose that B is total in \(F = G'\) for the topology \(\sigma(\mathrm{G}',\mathrm{G})\) ; if H is the subspace of \(F = G'\) generated by B, the hypothesis then implies that every bounded subset C of G is precompact and metrizable for \(\sigma(\mathrm{G},\mathrm{H})\) . Therefore every bounded sequence \((\mathbf{z}_{n})\) of points of G has a subsequence Cauchy for \(\sigma(\mathrm{G},\mathrm{H})\) ; show that if \((\mathbf{z}_{n})\) is a Cauchy sequence for \(\sigma(\mathrm{G},\mathrm{H})\) , and if one writes \(z_{n} \circ m = g_{z_{n}} \cdot \mu\) , the sequence of classes of the \(g_{z_{n}}\) in \(L^{1}(\mu)\) is a Cauchy sequence for the topology \(\sigma(\mathrm{L}^{1},\mathrm{L}^{\infty})\) ; finally, conclude a contradiction from this by using Exers. 15 and 16 of Ch. V, §5.)

\begin{enumerate}
\def\labelenumi{\arabic{enumi})}
\setcounter{enumi}{13}
\tightlist
\item
  \begin{enumerate}
  \def\labelenumii{\alph{enumii})}
  \tightlist
  \item
    Let G be a Banach space, \(F = G'\) its strong dual, g a mapping of T into F belonging to the space \(\Lambda_{G'}^{p}\) (§1, Exer. 16). Show that if p \textgreater{} 1, the measure \(g \cdot \mu\) is absolutely continuous with respect to \(\mu\) (for the strong topology of F).
  \end{enumerate}
\end{enumerate}

\begin{enumerate}
\def\labelenumi{\alph{enumi})}
\setcounter{enumi}{1}
\tightlist
\item
  Take for G the Banach space \(\mathrm{L}^{1}(\mathbf{N})\) , so that \(\mathrm{G}^{\prime}=\mathrm{L}^{\infty}(\mathbf{N})\) ; if f is the function defined in Exer. 7a) of §1 (regarded as taking its values in \(G^{\prime}\) ), then f is scalarly well integrable but the measure \(f\cdot\mu\) is not absolutely continuous with respect to \(\mu\) for the strong topology of \(G^{\prime}\) .
\end{enumerate}

¶ 15) a) Let f be a scalarly locally integrable mapping of T into a quasi-complete Hausdorff locally convex space F. Then \(m = f \cdot \mu\) is a vectorial measure with values in \(F'^{*}\) (for the topology \(\sigma(F'^{*}, F')\) ). Show that if this measure is absolutely continuous (for the topology of uniform convergence on the equicontinuous subsets of \(F'\) ), and if, moreover, f is \(\mu\) -measurable (for the original topology of F), then m is a vectorial measure with values in F. (Make use of Prop. 8 of §1, No.~2.) If F is a Banach space and if, for some p \textgreater{} 1, \(\langle z', f \rangle\) belongs to \(\overline{\mathcal{L}}^{p}(\mu)\) for all \(z' \in F'\) , then the measure m is absolutely continuous (cf.~Exer. 14 a)).

\begin{enumerate}
\def\labelenumi{\alph{enumi})}
\setcounter{enumi}{1}
\tightlist
\item
  Let \(\mathrm{I} = [0,1]\) , \(f\) the function defined on \(\mathrm{I} \times \mathrm{I}\) (using the continuum hypothesis) such that, for every \(t \in \mathrm{I}\) , \(s \mapsto f(s,t)\) is the characteristic function of a countable set, and, for every \(s \in \mathrm{I}\) , \(t \mapsto f(s,t)\) is the characteristic function of the complement of a countable set (Ch. V, §8, Exer. 7 c)). Denote by \(\mathrm{I}_0\) the interval I equipped with the discrete topology, by F the Banach space \(\mathcal{B}(\mathrm{I}) = \mathrm{L}^\infty(\mathrm{I})\) of bounded functions on I; one can identify F with the space \(\mathcal{C}(\widetilde{\mathrm{I}}_0)\) of continuous functions on the Stone-Čech compactification of \(\mathrm{I}_0\) . For every \(s \in \mathrm{I}\) , denote by \(\mathbf{g}(s)\) the element \(t \mapsto f(s,t)\) of F; show that \(\mathbf{g}\) is scalarly integrable for the Lebesgue measure \(\mu\) on I; but \(\int \mathbf{g} d\mu\) does not belong to F. More precisely, \(\mathbf{g} \cdot \mu\) is a vectorial measure with values in \(\mathrm{F}''\) , equal to \(\mathbf{c}\mu\) , where \(\mathbf{c}\) is the constant vector identified with the characteristic function of \(\widetilde{\mathrm{I}}_0 - \mathrm{I}_0 = \mathrm{E}\) (cf.~Exer. 12 a)). (Decompose every positive measure on \(\widetilde{\mathrm{I}}_0\) as the sum of an atomic measure and a diffuse measure (Ch. V, §5, No.~10, Prop. 15), and observe that the diffuse measure is carried by E.) From this, deduce that there does not exist any \(\mu\) -measurable function \(\mathbf{h}\) with values in \(\mathbf{F}\) such that \(\mathbf{h} - \mathbf{g}\) is scalarly \(\mu\) -negligible (make use of \(a\) )).
\end{enumerate}

\begin{enumerate}
\def\labelenumi{\arabic{enumi})}
\setcounter{enumi}{15}
\item
  Consider the mappings \(u_{m}\) of T = {[}0,1{]} into a separable Hilbert space, defined in §1, Exer. 1, and set \(f_{k} = u_{1} + \cdots + u_{k}\) ; show that the vectorial measures \(f_{k} \cdot \mu\) converge uniformly on every bounded subset of \(\mathcal{C}(T)\) to a vectorial measure m with values in F. Moreover, m may be extended by continuity to the space \(\mathcal{L}^{2}(\mu)\) and, for every function \(f \in L^{2}\) , one has \(\int f dm \in F\) and \(|\int f dm| \leqslant N_{2}(f)\) in F; in particular, m is scalarly of base \(\mu\) and absolutely continuous with respect to \(\mu\) for the strong topology (Exer. 12 c)). However, m is not of base \(\nu\) for any positive measure \(\nu\) on T, and a fortiori (No.~5, Cor. 4 of Th. 1) is not majorizable. (Reduce to the case that \(\nu\) has base \(\mu\) , using Th. 3 of Ch. V, §5, No.~7.)
\item
  Let \(\mu\) be Lebesgue measure on \(T = [0,1]\) ; the mapping \(f \mapsto f \cdot \mu\) of \(\mathcal{C}(T)\) into \(\mathcal{M}(T)\) is continuous for the strong topology on \(\mathcal{M}(T)\) , hence is a vectorial measure \(\mathbf{m}\) with values in that Banach space. Show that \(\mathbf{m}\) is scalarly of base \(\mu\) and absolutely continuous with respect to \(\mu\) for the strong topology, but that \(\mathbf{m}\) does not have base \(\mu\) (for the strong topology on \(\mathcal{M}(T)\) ). (Observe that \(\mathbf{m}\) has base \(\mu\) for the vague topology \(\sigma(\mathcal{M}(T), \mathcal{C}(T))\) , and deduce from this that if one had \(\mathbf{m} = \mathbf{g} \cdot \mu\) , with \(\mathbf{g}\) scalarly \(\mu\) -integrable (for the strong topology on \(\mathcal{M}(T)\) ), one would necessarily have \(\mathbf{g}(t) = \varepsilon_t\) almost everywhere. Show that this leads to a contradiction, by observing that for every bounded numerical function \(\theta\) on \(T\) , the linear form \(\mathbf{z}' : \lambda \mapsto \sum_{t \in T} \theta(t) \lambda(\{t\})\)
\end{enumerate}

on \(\mathcal{M}(\mathrm{T})\) is continuous but that \(\langle \mathbf{g},\mathbf{z}'\rangle\) is not necessarily \(\mu\) -measurable.)

¶ 18) a) Let T be a locally compact space, \((K_{\alpha})\) a locally finite covering of T by compact sets. Let \(\mu\) be a positive measure on T, \(\mu_{\alpha}\) the measure induced by \(\mu\) on \(K_{\alpha}\) . Assume that each of the spaces \(L^{\infty}(K_{\alpha}, \mu_{\alpha})\) has the lifting property. Show that \(L^{\infty}(T, \mu)\) has the lifting property.

\begin{enumerate}
\def\labelenumi{\alph{enumi})}
\setcounter{enumi}{1}
\tightlist
\item
  Let K be a metrizable compact space, \(\mu\) a positive measure on K, \((\varpi_n)\) a fundamental sequence (Ch. IV, §5, Exer. 17) of finite partitions of K into integrable sets. For every integrable subset A of K and every element \(\widetilde{f} \in \mathrm{L}^1(\mu)\) , set \(\lambda_A(\widetilde{f}) = 0\) if \(\mu(A) = 0\) , and \(\lambda_A(\widetilde{f}) = (\mu(A))^{-1} \int_{A} f d\mu\) otherwise; for every \(n\) , set \(\rho_n(\widetilde{f}) = \sum_k \lambda_{A_k}(\widetilde{f}) \varphi_{A_k}\) , if
\end{enumerate}

\(\varpi_{n} = (\mathrm{A}_{k})\) . Show that the sequence \((\rho_{n}(\tilde{f}))\) of \(\mu\) -integrable functions converges almost everywhere to \(f\) . (Consider first the case that \(f\) is a continuous function. Given an arbitrary number \(a > 0\) , show next that the union B of all the sets A belonging to at least one of the partitions \(\varpi_{n}\) and such that \(a \cdot \mu(A) \leqslant \int_{A} f d\mu\) , is measurable and is such that \(\mu(B) \leqslant a^{-1} \int f d\mu\) . Then approximate \(f\) in \(\mathcal{L}^{1}\) by a sequence of continuous functions.)

\begin{enumerate}
\def\labelenumi{\alph{enumi})}
\setcounter{enumi}{2}
\item
  Show that every metrizable compact space has the lifting property. (Let \(\mathfrak{U}\) be an ultrafilter on \(\mathbf{N}\) finer than the Fréchet filter; show that \(\rho(\widetilde{f}) = \lim_{\mathfrak{U}} \rho_n(\widetilde{f})\) is a lifting of \(\mathrm{L}^{\infty}\) .)
\item
  Deduce from \(a\) ) and \(c\) ) that every metrizable locally compact space has (for any positive measure whatever) the lifting property.
\end{enumerate}

\begin{enumerate}
\def\labelenumi{\arabic{enumi})}
\setcounter{enumi}{18}
\tightlist
\item
  Let \(\mu\) be a measure on T such that the Banach space \(L^1 (\mu)\) is separable. Show that every continuous linear mapping of \(L^1 (\mu)\) into the strong dual \(F'\) of an arbitrary Banach space \(F\) may be obtained by passage to the quotient starting from a mapping of the form \(g\mapsto \int g\mathbf{f}d\mu\) , where \(\mathbf{f}\in \mathcal{L}_{\mathrm{F}_s'}^\infty\) . (Reduce to the Dunford-Pettis theorem, using Exer. 19 of TVS, IV, §2.)
\end{enumerate}

¶ 20) Let F be a separable Fréchet space, F' its dual, μ a positive measure on T, m a measure on T with values in the weak dual \(F'_s\), scalarly of base μ. For m to have base \(\mu\) , it is necessary and sufficient that the following condition be fulfilled: for every \(\varepsilon > 0\) and every compact subset \(K\) of \(T\) , there exists a compact subset \(K_1 \subset K\) such that \(\mu(K - K_1) \leqslant \varepsilon\) , and such that the image under \(m\) of the set of \(\mu\) -measurable functions \(g\) that are bounded, have support contained in \(K_1\) and satisfy \(\int |g| d\mu \leqslant 1\) , is an equicontinuous subset of \(F'\) . (Recall that \(\int g dm \in F'\) for every \(\mu\) -measurable function \(g\) that is bounded and has compact support; cf.~Cor. 2 of Prop. 3. To show that the condition is necessary, use Prop. 13 of §1, No.~5 and Prop. 5 of §1, No.~2. To see that the condition is sufficient, consider first the case that \(T\) is compact; applying Cor. 3 of Th. 1, No.~5, first show that there exist a partition of \(T\) into a negligible set \(N\) and a sequence ( \(K_n\) ) of compact sets, and a measurable mapping \(g\) of \(T\) into \(F'_s\) , such that \(\int f dm = \int g f d\mu\) for every function \(f \in \mathcal{L}^\infty(\mu)\) that is zero on the complement of the union of a finite number of the sets \(K_n\) ; to prove that \(m = g \cdot \mu\) , use Exers. 16 b) and 20 a) of Ch. V, §5. Finally, to pass to the case that \(T\) is any locally compact space, use Prop. 14 of Ch. IV, §5, No.~9.)

¶ 21) Let F be a Banach space, u a continuous linear form on the Banach space \(\mathrm{L}_{\mathrm{F}}^{p}(\mathrm{T},\mu)\) ; denote by q the conjugate exponent of p, by E the space of numerical step functions over the \(\mu\) -integrable sets. For \(f \in L^{p}\) and \(z \in F\) , one can write \(u(fz) = \langle z, m(f) \rangle\) , where m is a continuous linear mapping of \(L^{p}\) into \(F'\) such that \(|m(f)| \leqslant \|u\| \cdot N_{p}(f)\) .

\begin{enumerate}
\def\labelenumi{\alph{enumi})}
\tightlist
\item
  Let \((\mathrm{A}_i)_{1\leqslant i\leqslant n}\) be a finite sequence of \(\mu\) -integrable sets, pairwise disjoint and non-negligible. Show that
\end{enumerate}

\[
\sum_ {i} \left(| \mathbf {m} (\varphi_ {\mathrm{A} _ {i}}) | ^ {q} / (\mu (\mathrm{A} _ {i})) ^ {q - 1}\right) \leqslant \| u \| ^ {q}.
\]

(For every system \((\mathbf{a}_i)_{1\leqslant i\leqslant n}\) of vectors in F such that \(|\mathbf{a}_i| = 1\) , consider the linear form \((\xi_i)\mapsto u\big(\sum_i\xi_i\mathbf{a}_i\varphi_{\mathrm{A}_i}\big)\) on \(\mathbf{R}^n\) , and apply Th. 4 of Ch. V, §5, No.~8 to a measure with finite support.) From this, deduce that if \(\mathrm{A} = \bigcup_{i}\mathrm{A}_{i}\) , then \(\sum_{i}|\mathbf{m}(\varphi_{\mathrm{A}_i})|\leqslant \| u\| (\mu (\mathrm{A}))^{1 / p}\) (apply Minkowski's inequality).

\begin{enumerate}
\def\labelenumi{\alph{enumi})}
\setcounter{enumi}{1}
\tightlist
\item
  For every positive function \(f \in \mathcal{E}\) , set \(\nu(f) = \sup\left(\sum_{i} |\mathbf{m}(f\varphi_{\mathrm{A}_i})|\right)\) , where \((\mathrm{A}_i)\)
\end{enumerate}

runs over the set of finite sequences of pairwise disjoint \(\mu\) -integrable sets. Show that \(\nu\) is the restriction to \(\mathcal{E}\) of a positive measure with base \(\mu\) (again denoted \(\nu\) ) on T, such that \(|\mathbf{m}(f)| \leqslant \nu(f)\) for every positive function \(f \in \mathcal{E}\) (make use of \(a\) )).

\begin{enumerate}
\def\labelenumi{\alph{enumi})}
\setcounter{enumi}{2}
\tightlist
\item
  Show that if F is separable, there exists a mapping g of T into F', scalarly μ-measurable, such that \(u(\widetilde{\mathbf{f}})=\int\langle\mathbf{f},\mathbf{g}\rangle d\mu\) for every function \(f\in L_{F}^{p}\) , and that \(\|u\|=N_{q}(\mathbf{g})\) . (To establish the last equality, use Prop. 13 of §1, No.~5 and Exer. 13 of §1.) d) Deduce from c) that if F is a reflexive separable Banach space, the space \(L_{F}^{p}\) is reflexive for \(1<p<+\infty\) (apply Prop. 12 of §1, No.~5).
\end{enumerate}

¶ 22) Let F be a Hausdorff locally convex space, S the set of subsets of F that are convex, balanced, and compact for \(\sigma(F, F')\) , and \(S'\) the set of subsets of \(F'\) that are convex, equicontinuous, balanced, and compact for \(\sigma(F', F'')\) .

\begin{enumerate}
\def\labelenumi{\alph{enumi})}
\item
  Suppose that every set in \(\mathfrak{S}\) is precompact for the \(\mathfrak{S}'\) -topology. Show that every continuous linear mapping \(u\) of \(F\) into a quasi-complete, Hausdorff locally convex space \(G\) , that transforms the bounded subsets of \(F\) into subsets relatively compact for \(\sigma(G, G')\) , transforms every set belonging to \(\mathfrak{S}\) into a subset of \(G\) that is relatively compact for the original topology. (Use Exer. 4 of TVS, IV, §1 and Exer. 11 above.) Converse (consider the canonical mapping of \(F\) into its completion for the \(\mathfrak{S}'\) -topology).
\item
  Suppose that F is a metrizable space, or a strict direct limit of metrizable spaces. Show that the condition of a) is satisfied if every Cauchy sequence for \(\sigma(\mathrm{F},\mathrm{F}^{\prime})\) is a Cauchy sequence for the \(S^{\prime}\) -topology (make use of Šmulian's theorem (TVS, IV, §5, Exer. 2 c))).
\item
  Suppose F is infra-barreled (TVS, III, §4, Exer. 7) and denote by \(\mathfrak{S}''\) the set of subsets of \(\mathrm{F}''\) that are convex, equicontinuous, balanced, and compact for \(\sigma(\mathrm{F}''\) , \(\mathrm{F}'''\) ). Show that if every set in \(\mathfrak{S}'\) is precompact for the \(\mathfrak{S}''\) -topology, then every set in \(\mathfrak{S}\) is precompact for the \(\mathfrak{S}'\) -topology. (Use Exer. 4 of TVS, IV, §1, and observe that the original topology of F is induced by the strong topology on \(\mathrm{F}''\) .)
\end{enumerate}

¶ 23) Let T be a locally compact space, \(\mu\) a positive measure on T, and F one of the three Banach spaces \(\overline{\mathcal{K}(T)}\) , \(L^1(\mu)\) , \(L^\infty(\mu)\) . Let u be a continuous linear mapping of F into a quasi-complete space G, that transforms the unit ball of F into a subset of G that is relatively compact for \(\sigma(G, G')\) . Show that u transforms every subset of F relatively compact for \(\sigma(F, F')\) into a subset of G relatively compact for the original topology. (Apply Exer. 22 c) to reduce the case \(F = L^1(\mu)\) to the case \(F = L^\infty(\mu)\) ; use Exer. 13 of Ch. II, §1 to reduce the case \(L^\infty(\mu)\) to the case \(F = \mathcal{C}(S)\) , where S is a suitable compact space. If \(F = \overline{\mathcal{K}(T)}\) , apply Exer. 22 b), then use Exers. 22 b) and 15 of Ch. V, §5; note that a Cauchy sequence for \(\sigma(F, F')\) converges pointwise on T, and a fortiori converges in measure for every bounded measure on T.)

¶ 24) Let \(\mu\) be a positive measure on \(T\) , \(u\) a linear mapping of \(L^1(\mu)\) into a Banach space \(F\) , transforming the unit ball of \(L^1(\mu)\) into a subset of \(F\) relatively compact for \(\sigma(F, F')\) . Show that there exists a \(\mu\) -measurable mapping \(f\) of \(T\) into \(F\) such that \(|\mathbf{f}(t)| \leqslant \|u\|\) for all \(t \in T\) and such that

\[
u (\widetilde {g}) = \int \mathbf {f} g d \mu
\]

for every function \(g \in \overline{\mathcal{L}}^1(\mu)\) (Dunford-Pettis-Phillips theorem). (First reduce to the case that T is compact, with the help of Prop. 14 of Ch. IV, §5, No.~9. The unit ball B of \(L^\infty(\mu) \subset L^1(\mu)\) is then relatively compact for \(\sigma(L^1, L^\infty)\) (Ch. V, §5, Exer. 15); using Exer. 23, show that the closed linear subspace of F generated by \(u(L^1)\) is separable. One can therefore reduce to the case that F is separable and \(\overline{u(L^1)} = F\) . The closure A in F of the image under \(u\) of the unit ball of \(L^1\) is then compact for \(\sigma(F, F')\) (Ch. IV, §7, Exer. 10); show that it is metrizable for this topology (TVS, III, §3, No.~4, Cor. 2 of Prop. 6). Observe that A may then be identified with the unit ball of the dual of the normed space G obtained by equipping \(F'\) with the norm equal to the gauge of \(A^\circ\) . Show that G is separable, and thus reduce to applying the Dunford-Pettis theorem (Cor. 2 of Th. 1).

¶ 25) Let \(\mu\) be a positive measure on T.

\begin{enumerate}
\def\labelenumi{\alph{enumi})}
\item
  Let \(\mathbf{f}\) be a scalarly \(\mu\) -measurable mapping of \(\mathrm{T}\) into a Banach space \(\mathrm{F}\) , such that \(\mathbf{f}(\mathrm{T})\) is relatively compact for \(\sigma(\mathrm{F},\mathrm{F}')\) . Show that there exists a \(\mu\) -measurable mapping \(\mathbf{g}\) of \(\mathrm{T}\) into \(\mathrm{F}\) such that \(\mathbf{f} - \mathbf{g}\) is scalarly locally negligible. (Consider the mapping \(\widetilde{h} \mapsto \int \mathbf{f}h d\mu\) of \(\mathrm{L}^1(\mu)\) into \(\mathrm{F}\) , and apply Exer. 24 to it, also using Exer. 10 of Ch. IV, §7.)
\item
  Let \(\mathbf{f}\) be a scalarly \(\mu\) -measurable mapping of T into a reflexive Fréchet space F. Show that there exists a \(\mu\) -measurable mapping \(\mathbf{g}\) of T into F such that \(\mathbf{f} - \mathbf{g}\) is scalarly locally negligible (cf.~Exer. 15 b)). (Reduce to the case that T is compact, using Prop. 14 of Ch. IV, §5, No.~9. Next, apply Exer. 23 c) of §1 to reduce to the case that \(\mathbf{f}(T)\) is bounded in F. Then embed F in a countable product of Banach spaces, and use a).)
\item
  Let \(\mathbf{f}\) be a mapping of \(\mathrm{T}\) into a Fréchet space \(\mathrm{F}\) , \(\mu\) -measurable for the topology \(\sigma(\mathrm{F}, \mathrm{F}')\) . Show that \(\mathbf{f}\) is \(\mu\) -measurable for the original topology of \(\mathrm{F}\) . (Reduce to the case that \(\mathrm{T}\) is compact and \(\mathbf{f}\) is continuous for the topology \(\sigma(\mathrm{F}, \mathrm{F}')\) ; conclude as in \(b\) .)
\end{enumerate}

¶ 26) Let S, T be two compact spaces, f a finite numerical function defined on S × T.

\begin{enumerate}
\def\labelenumi{\alph{enumi})}
\item
  In order that the mapping \(s \mapsto f(s, \cdot)\) of S into \(\mathbf{R}^{\mathrm{T}}\) be a continuous mapping of S into the space \(\mathcal{C}(\mathrm{T})\) equipped with the topology \(\sigma(\mathcal{C}(\mathrm{T}), \mathcal{M}(\mathrm{T}))\) , it is necessary and sufficient that \(f\) be bounded and that each of the partial mappings \(f(s, \cdot), f(\cdot, t)\) ( \(s \in \mathrm{S}, t \in \mathrm{T}\) ) be continuous. (To see that the condition is sufficient, first show that the image M of S under \(s \mapsto f(s, \cdot)\) is compact for \(\sigma(\mathcal{C}(\mathrm{T}), \mathcal{M}(\mathrm{T}))\) . For this, observe that \(s \mapsto f(s, \cdot)\) is continuous when \(\mathcal{C}(\mathrm{T})\) is equipped with the topology of pointwise convergence on T; use Eberlein's theorem (TVS, IV, §5, No.~3, Th. 1) to reduce to proving that every sequence \((f(s_n, \cdot))\) has a cluster point in \(\mathcal{C}(\mathrm{T})\) for \(\sigma(\mathcal{C}(\mathrm{T}), \mathcal{M}(\mathrm{T}))\) . Reduce to the case that T is metrizable, by considering a suitable quotient space of T (cf.~Exer. 10 a)), and observe that, on a subset of \(\mathcal{C}(\mathrm{T})\) relatively compact for the topology of pointwise convergence on T, this topology is identical to the topology of pointwise convergence on a dense subset of T. Obtain thus a subsequence of \((f(s_n, \cdot))\) that converges pointwise on T, and apply Lebesgue's theorem. Finally, observe that on M, the topology \(\sigma(\mathcal{C}(\mathrm{T}), \mathcal{M}(\mathrm{T}))\) and the topology of pointwise convergence are identical.)
\item
  Assume that each of the partial mappings \(f(s, \cdot), f(\cdot, t)\) is continuous ( \(s \in S\) , \(t \in T\) ). Show that for every positive measure \(\mu\) on \(S\) and every \(\varepsilon > 0\) , there exists a compact subset \(K \subset S\) such that \(\mu(S - K) \leqslant \varepsilon\) and the restriction of \(f\) to \(K \times T\) is continuous. (Reduce to the case that \(f\) is bounded; apply a) and Exer. 25 c).
\item
  With the same hypotheses as in b), show that f is measurable for every measure \(\nu\) on S × T. (Apply b) to the image of \(\nu\) under the projection of S × T onto S.)
\end{enumerate}

\begin{enumerate}
\def\labelenumi{\arabic{enumi})}
\setcounter{enumi}{26}
\tightlist
\item
  Let \(\mathbf{m}\) be a vectorial measure on \(\mathbf{T}\) with values in a Hausdorff locally convex space \(\mathbf{F}\) .
\end{enumerate}

\begin{enumerate}
\def\labelenumi{\alph{enumi})}
\item
  Show that if the support of \(\mathbf{m}\) is finite, then \(\mathbf{m} = \sum_{i=1}^{n} \mathbf{c}_i \varepsilon_{a_i}\) , where the \(\mathbf{c}_i \in \mathbf{F}\) .
\item
  If \(\mathbf{F}\) is a Banach space, and \(\mathbf{m}\) is continuous for the topology of compact convergence, show that \(\mathbf{m}\) has compact support.
\item
  Give an example of a measure with values in \(\mathbf{R}^{\mathbf{N}}\) , with noncompact support, that is continuous for the topology of compact convergence.
\end{enumerate}

\setcounter{exercisegroup}{3}
\edef\CurrentExerciseGroup{\arabic{exercisegroup}}
